\section{Auswertung der Laufversuche}
\label{sec:auswertung_laufversuche}

Die Laufversuche sind in einer Arbeitsmappe des Prüfstands festgehalten, je Lauf eine Zeile. Ihre Spalten nennen die Laufnummer, das Datum, den Prüfstand KV3 oder KV4, den Werkstoff, die ME-Nummern von Ritzel und Rad, die Öltemperatur, die Zieltemperatur des Zahns sowie die Versagensart. Es folgen das Zieldrehmoment am Ritzel, die Drehmomente nach dem Start sowie nach $15$\,min und $60$\,min und der Vermerk des Nachspannens. Den Schluss bilden das Raddrehmoment, die Drehzahl, die Lastwechselzahl und eine Bemerkung mit der Nachspann-Historie. Ein Einleseskript überführt die Mappe in eine geprüfte Tabelle. Eine Spaltenkarte ordnet die Überschriften festen Feldern zu. Zahlen mit Tausendertrennzeichen, Drehmomentbereiche wie $11{,}6$ bis $12{,}5$\,Nm und Texteinträge in Zahlenspalten werden in Zahlenwerte, in Bereiche und in Vermerke zerlegt. Das Raddrehmoment ergibt sich aus dem Ritzelmoment mit dem Zähnezahlverhältnis $52/51$. Gerundet auf die Stufen $8$\,Nm, $12$\,Nm und $16$\,Nm dient es als Stufenschlüssel, mit dem die Läufe den Lasthorizonten zugeordnet werden. Die Vergleichsläufe erhalten auch die Stufen $10$\,Nm, $14$\,Nm und $18$\,Nm. Doppelt vergebene Laufnummern, leere Platzhalterzeilen, eine abweichend geschriebene Laufnummer und Zeilen ohne Laufnummer, die zum vorangehenden Lauf gehören, werden erkannt. Sie stehen in einer Anomalieliste, die vor der Auswertung durchgesehen und bestätigt wird. Aus dem Freitext der Bemerkung wird die Nachspann-Historie gelesen, also Lastwechselzahl, gemessenes Moment vor dem Nachspannen, eingestelltes Moment danach und die Temperatur. Die Spanne der gemessenen Momente zwischen den Nachspannpunkten bleibt je Lauf als Bereich in der Tabelle, Lastgröße der Regression ist der Stufenwert. Der Zähler des Prüfstands zählt die Umdrehungen der Ritzelwelle. Die Drehzahlen $3000$ und $2250$\,1/min gelten an der Radwelle. Ein Zahn des Rades erfährt je Zählerstand $51/52$ Lastwechsel. Die Tabellen führen die Zählerstände. Die Umrechnung verschiebt $\lg N_{\mathrm{L}}$ um $0{,}008$ und bleibt unter der Auflösung der Auswertung.

Drei Versagensdefinitionen werden getrennt geführt und mit eigenen Zeichen benannt, $N_{\mathrm{A}}$ für den ersten sichtbaren Anriss, $N_{\mathrm{K}}$ für das Prüfstandskriterium und $N_{\mathrm{B}}$ für den Zahnbruch oder die Abschaltung beim Versagensgeräusch. Für jeden Lauf gilt $N_{\mathrm{A}} \le N_{\mathrm{K}} \le N_{\mathrm{B}}$. Das Prüfstandskriterium ist der Abfall des Drehmoments um mehr als $1$\,Nm über 100\,000 Lastwechsel, den Fiederling für Zahnfußbrüche an POM-Rädern eingeführt hat, weil das schleichende Risswachstum keinen eindeutigen Ausfallzeitpunkt zulässt \cite[Abschnitt~6.1, S.~65]{Fie23}. Es wird hier unter Beobachtung übernommen, lässt Anrisse zu und liefert die Lastwechselzahl am Ende des Laufs. Der Totalausfall ist der Verlust der Momentenübertragung, bei Hasl der augenblickliche Abfall des gemessenen Lastmoments \cite[S.~91]{Has18}. Er wird vermerkt, wenn er eingetreten ist. Für den Nachweis nach Stommel et al. gilt als Versagen die Beanspruchung, die zu einem ersten makroskopischen Anriss führt \cite[S.~204]{Sto18}. Dafür werden aus der Bemerkung die Sichtkontrollen mit ihrem Zählerstand gelesen, und der Anriss wird als Intervall zwischen der letzten Kontrolle ohne und der ersten mit sichtbaren Anrissen geführt. Die versagten Zähne jedes Rades werden nach dem Lauf in einer Sichtprüfung aufgenommen. Zahn 1 ist an den Teilen markiert. Der Konturscan misst die Zähne 1, 18 und 36 mit den Flanken ihrer Nachbarn, und die in der Bemerkungsspalte genannten Messzähne 1 bis 3 werden diesen drei Zähnen zugeordnet.

% TK erledigt (Kriterium nach Fiederling unter Beobachtung; Fenster gestrichen; Messzähne 1/2/3 = Zahn 1/18/36 als Annahme, bitte bestätigen): Nicht formulieren das Kriterium des Lehrstuhls. Das ist von Fiederling und ich habe es mit unter Beobachtung aufgenommen. außerdem wieder das Fenster rausnehmen, das ist a priori annahme und gehört an der stelle immer noch nicht rein.

Die Wöhlerlinie wird im Zeitfestigkeitsbereich aus den drei Horizonten $8$\,Nm, $12$\,Nm und $16$\,Nm am Rad mit wenigen Wiederholungen je Horizont bestimmt. Diese Belegung ist die Ausgangslage der Arbeit und wird bei jeder Zahl mitgenannt. Die Belegung mit zwei, drei und vier Rädern auf drei Horizonten ist eine Mischform aus Horizonten- und Perlenschnurverfahren. DIN 50100 gleicht Ergebnisse auf mehr als zwei Horizonten gemeinsam durch eine Regression aus, wie beim Perlenschnurverfahren mit einer Probe je Horizont \cite[Abschnitt~10.2.3, S.~57]{DIN50100}. Angepasst werden zwei Geraden. Die halblogarithmische Gerade mit linearer Lastachse über $\lg N_{\mathrm{L}}$ folgt der Darstellung der Zeitschwellfestigkeit in VDI 2736 Blatt 2 und der Auswertung bei Hasl \cite[S.~107]{Has18}:
\begin{equation}
    \lg N_{\mathrm{L},50\%} = a_{\mathrm{c}} - b_{\mathrm{c}}\,T_{2} ,
    \label{eq:halblog_gerade}
\end{equation}
mit den Konstanten $a_{\mathrm{c}}$ und $b_{\mathrm{c}}$ der Geradengleichung \cite[S.~XIII und 110]{Has18}. Die doppeltlogarithmische Gerade nach \eqref{eq:zeitfestigkeitsgerade} liefert den Neigungsexponenten $k$, der an die Stelle der synthetischen Wöhlerlinie des Nachweises in \secref{sec:auswertung_beanspruchung} treten kann. Gepoolt wird die Streuung von $\lg N_{\mathrm{L}}$ über alle Horizonte. Dazu werden die Ergebnisse nach \eqref{eq:perlenschnur_streuung} parallel zur Ausgleichsgeraden auf einen gemeinsamen Horizont verschoben. Die Standardabweichung wird mit zwei verlorenen Freiheitsgraden gebildet und für den kleinen Stichprobenumfang korrigiert \cite[Abschnitt~10.2.2, S.~56 bis 57]{DIN50100}. Die Norm stuft Stichproben unter zehn Proben für die Streuung als unzuverlässig ein \cite[S.~30]{DIN50100}. Deshalb wird die Standardabweichung nach unten auf $0{,}10$ begrenzt, wie Hasl es für kleine Ausfallwahrscheinlichkeiten durchführt \cite[S.~118]{Has18}. Der Wert wird gegen den üblichen Bereich von $0{,}10$ bis $0{,}30$ gestellt \cite[S.~30]{DIN50100}. Aus Mittelwert und korrigierter Standardabweichung folgen nach \eqref{eq:quantil} die Linien für $50$\,\%, für $10$\,\% und für $1$\,\% Ausfallwahrscheinlichkeit, parallel zur Ausgleichsgeraden \cite[S.~63]{DIN50100}. Das Quantil für $10$\,\% ist $u = -1{,}282$ nach Tabelle 11 der Norm. Für $1$\,\% wird $u = -2{,}33$ gesetzt, wie Hasl es verwendet \cite[S.~9]{Has18}. Diese Linie wird als Extrapolation über den von der Norm gedeckten Bereich hinaus gekennzeichnet \cite[S.~29]{DIN50100}. Daneben wird die Umrechnung nach VDI 2736 Blatt 4 mit \eqref{eq:vdi2736_10prozent} je Horizont ausgeführt, damit der Wert für $10$\,\% nach der Richtlinie neben dem nach der Norm steht.

% TK erledigt (Mischform aus Horizonten- und Perlenschnurverfahren erklärt): Erklär mir das nochmal mit dem Perlschnurverfahren. Was ist damit gemeint, Ergebnisse auf zwei Horizonten. Wir haben doch immer nur ein Ergebnis auf einem Horizont und niemals bei 12 Nm und 8 Nm gleichzeitig? Oder was meinst du?

Als Kontrolle wird die Weibull-Verteilung auf denselben Daten geprüft. Dazu werden die Lastspielzahlen der Einzelversuche auf den Funktionswert der Ausgleichsgeraden ihres Horizonts bezogen, wie Hasl es für Versuche ohne mehrfach belegte Lastniveaus vorschlägt \cite[S.~109 bis 110]{Has18}. Die bezogenen Werte aller Horizonte werden gemeinsam nach ihrer Größe geordnet. Jedem Wert wird eine Ausfallwahrscheinlichkeit nach Rossow mit \eqref{eq:rossow} zugeordnet, zitiert bei Fiederling \cite[S.~34]{Fie23}. \Note{PH: Erstquelle Rossow 1964 (Qualitätskontrolle 9, Heft 12, S. 146 bis 147) beschaffen und direkt zitieren.} Die Punkte werden in ein Gauß-Netz über $\lg N_{\mathrm{L}}$ und in ein Weibull-Netz eingetragen. Gewählt wird für PA46-GF30 die Verteilung, deren Gerade die Punkte besser trifft, und die Wahl wird mit der Einschränkung der kleinen Stichprobe benannt.

% TK erledigt (Illenberger-Vergleich gestrichen, Rossow als Sekundärzitat mit Platzhalter für die Erstquelle): Inwieweit dienen die Werte von Illenberger als Vergleich? Und Rossow bitte wieder die Originalquelle und nicht Fiedler zitieren.

Durchläufer liegen nicht vor. Eine Verlängerung der Geraden über den belegten Bereich hinaus unterbleibt \cite[S.~16]{DIN50100}.

% gestrichen (TK kap3_8:25): Absatz zur Behandlung von Durchläufern.

% TK erledigt: Das Verfahern kannst du rausnehmen, da wir keine Durchläufer am Stichtag haben. Entsrpechend Andere Passagen, die darauf aufbauen ändern

Die Wöhlerbilder tragen die Lastspielzahl über drei Abszissen auf. Als erste dient das Raddrehmoment $T_{2}$ der Stufe, als zweite die Zahnfußspannung $\sigma_{\mathrm{F}}$ dieser Stufe je Berechnungsverfahren aus \secref{sec:analytik}. Dritte Abszisse ist die größte Hauptnormalspannung $\sigma_{1}$ des FE-Modells am maßgebenden Fußknoten aus \secref{sec:auswertung_beanspruchung}. Weil jede Spannung eine eindeutige Funktion des Drehmoments ist, bleiben die Lastspielzahlen und die Regression in $\lg N_{\mathrm{L}}$ unverändert. Nur die Lastachse wird punktweise umgerechnet. Zusätzlich werden die Zeitschwellfestigkeitswerte für PA66 nach VDI 2736 Blatt 2 eingetragen, die Hasl zur Beschreibung von PA46 heranzieht \cite[S.~118]{Has18}.

Neben PA46-GF30 werden allein die Läufe mit POM aus dem Jahr 2023 eingetragen \cite[S.~43 und 65]{Fie23}, nach \secref{sec:pruefstand} bei $73$ bis $76$\,°C Öltemperatur und $2250$\,1/min. Prüfstand, Ritzel, Drehzahl und Jahr stehen in der Legende, und die Lastwechselzahlen werden gegen die Protokolle der Masterarbeit geprüft. Ein gemeinsames Wöhlerbild aller Werkstoffe, Jahre und Prüfstände wird nicht gezeichnet, weil sich zwischen den Jahren Prüfstand, Ritzel, Drehzahl und Ölversorgung unterscheiden und ein Unterschied der Lastwechselzahlen keinem dieser Einflüsse zuzuordnen wäre. Verglichen werden die Lage und die Neigung der Geraden von PA46-GF30 und POM je Horizont, mit den genannten Unterschieden der Prüfbedingungen als Vorbehalt.

% TK erledigt (nur POM 2023, kein gemeinsames Bild aller Werkstoffe, Vergleich je Horizont mit Vorbehalt): Das vorgehen ist mir zu grob und habe ich weiter oben schon beschrieben, was gemacht werden soll. DAs ist wie der Vergelich von Äpfel und Birnen, wenn einfach alles innerhalb von einem Wöhlerdiagramm eingetragen wird. Bitte nochmal mit bereits vorhandenen Anmerkungen gegenprüfen und Auswertung verfeiner, bei Unklarheiten oder wenn du nicht weißt, was du machen sollst mich fragen!

Der Vorversuch zur Temperatur am Rad 97351 belegt eine Matrix aus den Raddrehmomenten $8$\,Nm, $12$\,Nm und $16$\,Nm und den Drehzahlen $2250$\,1/min und $3000$\,1/min am Rad. Je Zelle wurde nach zwei Stunden die Zahnmassentemperatur mit der Wärmebildkamera bei einem Emissionsgrad von $1$ abgelesen. Dazu steht die Öltemperatur an der Messstelle vor der Prüfverzahnung. Die sechs Werte werden als Balken je Drehmoment und Drehzahl aufgetragen und dem Bezug von $80$\,\textdegree C gegenübergestellt, auf den die Festigkeitswerte und die Werkstoffkarten bezogen sind. VDI 2736 Blatt 4 verlangt den werkstoffabhängigen Emissionskoeffizienten aus Referenzmessungen und lässt bei Ölschmierung die Öltemperatur als Näherung der Zahnfußtemperatur zu \cite[Abschnitt~8.2, S.~16 bis 17]{VDI2736-4}. Der Emissionsgrad wurde in einem Vorversuch der Forschungsstelle mit einem Messstreifen bestimmt und zu $1$ ermittelt. Mit diesem Wert sind die Aufnahmen ausgewertet. Die Aufnahmen stehen als Beispiel in \secref{sec:temperatur} und vollständig im Anhang.

% Herkunft:
% Absatz 1 (Mappe, Einlesen): 70_input/72_Kleingetriebepruefstand_KV3/kst-E_Laufversuche_KV3.xlsx, Blatt Tabelle1
%   (103 Zeilen, 21 Spalten; Kopfzeile: A, Datum, Prüfstand, Werkstoff, ME-Nr. Ritzel, ME-Nr. Rad, Öltemperatur,
%   Ziel/Zahntemperatur, Versagensart, Ziel/Drehmoment/Ritzel, verspannt, T_0 nach Start, nach 15 min, T/60 min,
%   T_nachgespannt, Drehmoment Rad, Drehzahl, Lastwechsel, Bemerkung, Grübchenfläche); Stufenschlüssel 52/51 aus
%   Spalte Drehmoment Rad (7,9 → 8,05; 11,8 → 12,03; 15,7 → 16,01); Mischtypen: Lastwechsel als Text mit
%   Leerzeichen (Zeile 13 ff.), Bereiche "11,6-12,5" (Zeile 77), "nicht vermerkt" (Zeile 82), "-" (Zeile 23);
%   Doppelzeilen: kst-E_54 bis kst-E_62 leer (Zeilen 55 bis 63) und gefüllt (Zeilen 64 bis 76), kst-E_58 bis
%   kst-E_61 zweimal (Zeilen 68 bis 75, "Daten verlieren wegen Problem von Rechner"), "kst-E-67" (Zeile 87),
%   Zeilen ohne Laufnummer 78 bis 83 zu kst-E_63; Nachspann-Historie im Freitext (Zeilen 88 bis 95, Muster
%   "<LW> <Nm> --> nachverspannt auf <Nm> bei <°C>"); Roadmap L1 (implementation_roadmap.md Zeile 216) und Plan
%   Abschnitt 6 C4 (hi-kannst-du-mich-cozy-rabin.md Zeile 217). Gebaut 2026-10-18: measurement/test_rig.py und
%   measurements.py testruns (Spaltenkarte, Ereignisse, Anrissfenster, Anomalien); Bestätigungen des Nutzers in
%   data/measurement/test_runs_user.yaml. Drehzahlen: Spalte R trägt auf KV3 die Ritzeldrehzahl (3058 = 3000 am
%   Rad), auf KV4 die Raddrehzahl (Nutzer 2026-10-18); Wärmebildkamera statt Infrarotthermometer (Nutzer).
% Absatz 2 (Versagensdefinitionen): MA-0296.txt Druckseite 65 (6.1: Totalausfall "Loss-of-drive", schleichendes
%   Risswachstum bei Zahnfußbruch, Abfall > 1 Nm über 100 000 LW, Abbildung 53); Hasl2018.txt Druckseite 91 (7.2);
%   Stommel2018.txt Druckseite 204 (3.2: erster makroskopischer Anriss); Roadmap L7 (Zeile 223: drei Definitionen
%   getrennt); Fehlerfenster 36 bis 40 und Zahn 1 markiert: kap3_3_dokumentation.tex Absatz 1; Spalte Versagensart
%   "ZFB: Zahn 34" (Zeile 89), "Zahn 40 und 41" (90), "Zahn 31" (91), "Zahn 39" (92), "Zahn 40" (93);
%   Sichtkontrollen mit Zählerstand und "Anrisse" in den Bemerkungen der Zeilen 89, 93 bis 95; "Messzahn 3"
%   (Zeile 89) und "die drei markierten Zähne" (Zeile 94) ohne Zahnnummer, daher Platzhalter.
% Absatz 3 (Regression, Streuung, Linien): DIN50100.txt S. 57 (10.2.3), S. 56 bis 57 (Gl. (36), (38), (39)),
%   S. 30 (n < 10, 0,10 bis 0,30), S. 63 (Gl. (67), Tabelle 11, parallel), S. 29 (Extrapolation über 10 bis 90 %
%   unsicher); Hasl2018.txt Druckseite 107 (Darstellung linear über log), XIII (ac, bc), 110 (Gl. (51)), 118
%   (slog ≥ 0,1), 9 (Gl. (10) mit 2,33); VDI2736-4.txt S. 9 bis 10 (Gl. (10) bis (15)); Roadmap L2 (Zeile 217:
%   halblogarithmisch und Basquin, gepoolte Streuung, 50 → 10 und 1 %); Roadmap E8 (Zeile 209: lokale Wöhlerlinie
%   mit k statt A_W und k des Buchs, offen); Entscheidung 2026-10-11 (Zeile 233: wenige Wiederholungen als
%   Ausgangslage nennen). Die Gleichung eq:halblog_gerade ist Hasl Gl. (51) mit T_2 statt σ_F0 als Lastgröße.
% Absatz 4 (Weibull, Netze): Hasl2018.txt Druckseite 109 bis 110 (bezogene Lastspielzahl); MA-0296.txt Druckseite 34
%   (Gl. (16) Rossow); Illenberger2022.txt Druckseite 19 (Gl. (19)), 96 (Tabelle 7-2); Roadmap L2 (Verteilung
%   nach den Netzen wählen, mit Einschränkung nennen).
% Absatz 5 (Durchläufer): DIN50100.txt S. 16 (NG, keine Extrapolation über die belegten Horizonte), S. 54 (10.2.1);
%   Roadmap L3 (Zeile 218: Durchläufer als Pfeile), L7 (Zeile 223: zweiter 8-Nm-Lauf läuft noch); Stichtag offen
%   (kap6_laufversuche.tex Platzhalter). Die doppelte Anpassung (Normregel und alle Ausfälle) ist eine Festlegung
%   dieses Entwurfs, weil der 8-Nm-Horizont einen Ausfall und einen Durchläufer trägt.
% Absatz 6 (Abszissen): Roadmap L3 (Zeile 218: T_2, σ_F je Verfahren, FE σ1, VDI 2736 PA66-Band); Hasl2018.txt
%   Druckseite 118 (PA66-Werte für PA46); Verweise sec:analytik und sec:auswertung_beanspruchung.
% Absatz 7 (Vergleichswerkstoffe): Mappe Zeilen 2 bis 45 (KV4 2023, POM/PEEK, 2250 1/min, Ritzel 86481), Zeilen
%   51 bis 76 und 84 bis 87 (KV3 2026, POM/PEEK, 3058 1/min, Ritzel 95919; Zeilen 51 bis 53 Trockenlauf bei RT mit
%   1945 1/min); MA-0296.txt Druckseite 66 (6.2: POM Zahnfußbruch, PEEK Grübchen); Roadmap L5 (Zeile 220: Zeilen
%   2023 gegen MA-0296 prüfen, PEEK-Typ offen); Gliederung 6.4.
% Absatz 8 (Temperaturmatrix): Mappe Zeilen 77 bis 83 (kst-E_63, Rad 97351, Ritzel 95919, Zieldrehmoment Ritzel
%   11,8/7,9/7,8/15,7, Drehzahl 2294 und 3058, Bemerkung "nach 2h ... °C Zahnflankentemperatur mit Emissionsgrad 1
%   (nach Test von Liang)", Öltemperatur 75); Roadmap L4 (Zeile 219) und Plan C6 (Zeile 219); Gliederung 4.5
%   (Öl- und Kastentemperatur, Bezug 80 °C); VDI2736-4.txt S. 16 bis 17 (8.2, 8.2.2, 8.2.4). Die Kastentemperatur
%   steht in der Mappe nur in späteren Läufen (TA, TP in Zeilen 94 bis 95); welche Werte für die Matrix vorliegen,
%   zeigt die Auswertung.
% Formeln: eq:halblog_gerade = Hasl 2018 Gl. (51) S. 110 (N_L,50 % = 10^(a_c − σ_F0,50 % · b_c)), hier mit T_2;
%   verwiesen wird auf eq:zeitfestigkeitsgerade, eq:perlenschnur_streuung, eq:quantil, eq:vdi2736_10prozent und
%   eq:rossow aus kap2_6_experiment_statistik.tex.
% Bildwunsch: keiner (Bilder in Kapitel 6).
% Bibwunsch: keine neuen Schlüssel (DIN50100, VDI2736-2, VDI2736-4, Has18, Ill22, Fie23, Sto18 vorhanden).
% Nomenklatur: a_c, b_c, -, Konstanten der logarithmischen Geradengleichung, Hasl 2018 S. XIII (neu); T_2, N m,
%   Raddrehmoment (T mit Index 2 steht in Nomenklatur.tex); N_L, Lastwechselzahl (neu, siehe 2.6); sigma_F, MPa,
%   Zahnfußspannung (nicht in Nomenklatur.tex; prüfen, ob 2.5.1 oder 3.4 sie führt); sigma_1 steht; u, P_A, j, n,
%   k wie in 2.6 gemeldet; der Emissionsgrad bleibt ohne Zeichen.
% Offen: Keine BA-Übernahme in diesem Abschnitt. Nicht belegt und deshalb als Platzhalter: Welle des
%   Lastwechselzählers, Zuordnung der Messzähne, PEEK-Typ. Der Stichtag der Laufversuche fehlt (Platzhalter in
%   Kapitel 6). Die Festlegung, bei Horizonten mit Durchläufer zwei Fassungen der Geraden auszuweisen, und der
%   Mindestwert 0,10 der Standardabweichung sind Vorschläge dieses Entwurfs nach DIN 50100 S. 30 und 54 sowie
%   Hasl S. 118; der Nutzer bestätigt sie.
